% accessibility.tex
\section{Accessibility (WCAG 2.1)}
To ensure that NExAIE articles are accessible to all readers, including those using screen readers, we require all authors to follow these WCAG 2.1 standards.

\subsection{Implementing Accessible Figures}
Use the \texttt{\textbackslash nexaieFigure} command provided by the class. This command wraps the image in an accessibility layer (ActualText metadata).

\nexaieFigure{ht}{media/logo.pdf}{The official NExAIE logo demonstrating the accessible figure command.}{A circular gold and white logo featuring a large stylized H and a small red Greek letter Psi.}

\subsection{Accessible Tables}
Screen readers navigate tables linearly. To maintain accessibility:
\begin{itemize}
    \item \textbf{Avoid Merged Cells}: Avoid using \texttt{multirow} or \texttt{multicolumn} unless absolutely necessary.
    \item \textbf{Row Headers}: Ensure the first column clearly labels the data that follows in the row.
\end{itemize}

\subsection{Color and Contrast}
When creating figures:
\begin{itemize}
    \item \textbf{Non-Color Cues}: Do not use color as the only way to convey meaning. Use patterns, dashed lines, or different marker shapes in plots.
    \item \textbf{Contrast}: Ensure a contrast ratio of at least 4.5:1 for text against its background.
\end{itemize}

\subsection{Descriptive Hyperlinks}
Hyperlink text should be descriptive of the destination. 
\begin{itemize}
    \item \textbf{Poor}: Click \href{https://kahveci.pw/pedal}{here} to see the archive.
    \item \textbf{Good}: For further details, see the \href{https://kahveci.pw/pedal}{PEDAL Archive}.
\end{itemize}
